
%%%%%%%%%%%%%%%%%%%%%%%%%%%%%%
\section{Positional Notation}
%%%%%%%%%%%%%%%%%%%%%%%%%%%%%%

Positional notation is a method for representing numbers in which the position of each digit in a number determines its value. This is a foundational concept in arithmetic and computer science, as it enables efficient representation and manipulation of numbers. In positional notation, each position is associated with a power of a given base (or radix), allowing numbers to be represented compactly and with minimal symbols \cite{ross2008}.

%******************************************
\subsection{Positional Notation in Base 10}
%******************************************

In the decimal system, which is base 10, each digit in a number corresponds to a specific power of 10, depending on its position. For any integer \( N \) composed of \( n \) digits \( d_n, d_{n-1}, \ldots, d_1, d_0 \), we can write \( N \) as:
\[
N = d_n \cdot 10^n + d_{n-1} \cdot 10^{n-1} + \ldots + d_1 \cdot 10^1 + d_0 \cdot 10^0
\]
where each \( d_i \) is a digit contained in the set \(d \in D = \{0, 1, 2, 3, 4, 5, 6, 7, 8, 9\}
 \). 
 \vspace{\baselineskip}\par
 \( 10^n \) represents the highest place value, where \( d_n \) is the coefficient.
Each subsequent term decreases the power of 10 by one, moving from left to right in the number.

%*****************************************************
\subsection{Example of Positional Notation in Base 10}
%*****************************************************

To illustrate, let’s examine the number \( 357 \). We can expand this number based on its positional values as follows:
\[
357 = 3 \cdot 10^2 + 5 \cdot 10^1 + 7 \cdot 10^0
\]
Breaking down each part:
\begin{itemize}
    \item The leftmost digit, \( 3 \), is in the hundreds place. Its positional value is \( 3 \cdot 10^2 = 300 \).
    \item The middle digit, \( 5 \), is in the tens place, contributing \( 5 \cdot 10^1 = 50 \).
    \item The rightmost digit, \( 7 \), is in the units place, contributing \( 7 \cdot 10^0 = 7 \).
\end{itemize}
Adding these values gives:
\[
357 = 300 + 50 + 7
\]
which confirms that the positional notation accurately represents the number \( 357 \).

%*********************************************
\subsection{Importance of Positional Notation}
%*********************************************
The power of positional notation lies in its ability to represent large numbers compactly. Unlike systems that use tally marks or other non-positional symbols, positional notation can handle any magnitude of numbers with a finite set of symbols, scaling efficiently as numbers grow. For instance, to represent a number like \( 357 \), only three digits are required in base 10. In contrast, a tally mark system would require writing out 357 individual tally marks, which could be grouped for readability but remains far less efficient than positional notation.

%**********************************************************
\subsection{Positional Notation for Arithmetic Operations}
%**********************************************************
Positional notation also enables straightforward arithmetic operations such as addition, subtraction, multiplication, and division by aligning numbers by their place values. For example, adding two numbers in positional notation involves summing each pair of digits in the same place value and carrying over if the sum exceeds 9. This place-based structure makes arithmetic systematic and algorithmic, lending itself well to both manual and computational methods of calculation \cite{burton2011}.


%********************************
\subsection{Base/Radix Concepts}
%********************************
In mathematics and computer science, a base (also called a radix) is the number of unique symbols, including zero, used to represent numbers in a positional numeral system \cite{knuth1997,rosen2018}. The base of a number system is denoted by \( r \), and the unique symbols, or digits, are contained in the set \( D = \{0, 1, 2, \ldots, r-1\} \). In the decimal (base 10) system, for example, \( r = 10 \), and the set of digits \( D \) is \( \{0, 1, 2, 3, 4, 5, 6, 7, 8, 9\} \). However, in other bases, such as binary (base 2) or hexadecimal (base 16), the set \( D \) and the value of \( r \) would differ.

%*******************************************
\subsection{Generalized Positional Notation}
%*******************************************

For a number in any base \( r \), the representation follows a similar positional notation structure to base 10. A number \( N \) in base \( r \) with \( n+1 \) digits can be written as:
\[
N = d_n \cdot r^n + d_{n-1} \cdot r^{n-1} + \ldots + d_1 \cdot r^1 + d_0 \cdot r^0
\]
where each \( d_i \in D \) is an integer from \( 0 \) to \( r-1 \) and represents the coefficient of each power of \( r \). This means that each position in the number corresponds to a specific power of \( r \), with \( r^n \) representing the highest place value (determined by the leftmost digit, \( d_n \)) and each subsequent term decreasing by a power of \( r \).

%****************************************************
\subsection{Positional Notation for Different Bases}
%****************************************************
The decimal system (base 10) is just one example of a positional numeral system. Other bases follow the same general rules but use different sets of digits. For instance:
\begin{itemize}
    \item Binary (base 2): Uses \( D = \{0, 1\} \), with each position representing powers of 2. This system is fundamental in digital computing, as all binary values can be represented by combinations of 0 and 1. A number \( N \) in binary can be represented as:
    \[
    N = d_n \cdot 2^n + d_{n-1} \cdot 2^{n-1} + \ldots + d_1 \cdot 2^1 + d_0 \cdot 2^0
    \]

    \item Octal (base 8): Uses \( D = \{0, 1, 2, 3, 4, 5, 6, 7\} \), where each position represents powers of 8. A number \( N \) in octal can be represented as:
    \[
    N = d_n \cdot 8^n + d_{n-1} \cdot 8^{n-1} + \ldots + d_1 \cdot 8^1 + d_0 \cdot 8^0
    \]

    \item Hexadecimal (base 16): Uses \( D = \{0, 1, 2, 3, 4, 5, 6, 7, 8, 9, A, B, C, D, E, F\} \), where each position represents powers of 16. A number \( N \) in hexadecimal can be represented as:
    \[
    N = d_n \cdot 16^n + d_{n-1} \cdot 16^{n-1} + \ldots + d_1 \cdot 16^1 + d_0 \cdot 16^0
    \]

    \item Duodecimal (base 12): Uses \( D = \{0, 1, 2, 3, 4, 5, 6, 7, 8, 9, \text{A}, \text{B}\} \), where each position represents powers of 12. A number \( N \) in duodecimal can be represented as:
    \[
    N = d_n \cdot 12^n + d_{n-1} \cdot 12^{n-1} + \ldots + d_1 \cdot 12^1 + d_0 \cdot 12^0
    \]

    \item Sexagesimal (base 60): Uses \( D = \{0, 1, 2, \ldots, 59\} \), with each position representing powers of 60. A number \( N \) in sexagesimal can be represented as:
    \[
    N = d_n \cdot 60^n + d_{n-1} \cdot 60^{n-1} + \ldots + d_1 \cdot 60^1 + d_0 \cdot 60^0
    \]
\end{itemize}



Bases 2, 8, and 16 are often used in computing and digital electronics because of their implementation properties. Historically, base 60, or the sexagesimal system, was used by the ancient Babylonians for mathematical and astronomical calculations, and its influence persists today in the way we measure time (e.g., 60 seconds in a minute, 60 minutes in an hour) \cite{robson2008}. Base 12, known as the duodecimal system, has also been significant, with roots in ancient cultures where it was valued for its divisibility by multiple factors (2, 3, 4, and 6), making it useful for trade, measurements, and time-keeping systems. Traces of base 12 are still evident in modern units, such as a dozen (12 items) and the division of hours into sets of 12 on a clock face \cite{ifrah2000}.


